\subsection{动态参考点下的CPT目标}

令第 $t$ 期在当前市场环境、参考点 $r_t$ 和偏好约束 $c_t$ 下，由策略 $\pi_\theta$ 与环境交互得到的窗口累计净盈亏随机变量记为 $X_t$。定义相对收益随机变量
\begin{equation}
Y_t = X_t-r_t.
\end{equation}

于是，第 $t$ 期的目标函数定义为
\begin{equation}
J_t(\theta,r_t)=CPT_t(\theta,r_t),
\end{equation}
其中
\begin{equation}
\label{eq:CPT-objective}
\begin{aligned}
CPT_t(\theta,r_t)
&=
\int_0^\infty
w_g\!\left(
\mathbb P_\theta\!\bigl(v_g((X_t-r_t)_+) > z \,\big|\, s_t,c_t\bigr)
\right)\,dz \\
&\quad -
\int_0^\infty
w_l\!\left(
\mathbb P_\theta\!\bigl(v_l((r_t-X_t)_+) > z \,\big|\, s_t,c_t\bigr)
\right)\,dz .
\end{aligned}
\end{equation}

这里 $v_g(\cdot),v_l(\cdot)$ 分别为收益侧和损失侧的价值函数，$w_g(\cdot),w_l(\cdot)$ 分别为收益侧和损失侧的概率加权函数。对应地，第 $t$ 期的一阶驻点集合定义为
\begin{equation}
\mathcal S_t
=
\left\{
\theta\in\Theta:\nabla J_t(\theta,r_t)=0
\right\}.
\end{equation}

因此，本文研究的核心问题是：在动态参考点与非平稳真实市场环境下，构造一个基于 CPT-PG 思想的在线优化器，使得策略参数序列 $\{\theta_t\}$ 能够持续跟踪每期目标函数 $J_t(\theta,r_t)$ 的一阶驻点集合 $\mathcal S_t$。

\subsection{CPT策略梯度估计器}

在第 $t$ 日开盘前，优化器只使用滑动窗口
\begin{equation}
H_t=\{t-h,\dots,t-1\}
\end{equation}
中的历史数据进行估计。

参考 CPT-PG 的 Algorithm 1，在第 $t$ 期分别独立生成两组样本：

\begin{itemize}
\item 用于估计 CPT 权重的 $n_t$ 条独立收益样本：
\begin{equation}
\widetilde X_t^{(i)},\qquad i=1,\dots,n_t,
\end{equation}
并定义
\begin{equation}
\widetilde Y_t^{(i)}=\widetilde X_t^{(i)}-r_t,\qquad i=1,\dots,n_t.
\end{equation}

\item 用于估计策略得分的 $m_t$ 条独立轨迹样本：
\begin{equation}
X_t^{(j)},\qquad j=1,\dots,m_t,
\end{equation}
并定义
\begin{equation}
Y_t^{(j)}=X_t^{(j)}-r_t,\qquad j=1,\dots,m_t.
\end{equation}
第 $j$ 条轨迹对应的策略得分项记为
\begin{equation}
\sum_{k=1}^{h}
\nabla_\theta
\log \pi_{\theta_{t-1}}
\bigl(
a_{t-k}^{(j)}\mid s_{t-k}^{(j)},c_t
\bigr).
\end{equation}
\end{itemize}

基于用于估计 CPT 权重的样本 $\{\widetilde Y_t^{(i)}\}_{i=1}^{n_t}$，定义经验尾概率
\begin{equation}
\widehat p_t^g(y)
=
\frac{1}{n_t}
\sum_{i=1}^{n_t}
\mathbf 1\{\widetilde Y_t^{(i)}\ge y\},
\qquad y\ge 0,
\end{equation}
以及
\begin{equation}
\widehat p_t^l(y)
=
\frac{1}{n_t}
\sum_{i=1}^{n_t}
\mathbf 1\{\widetilde Y_t^{(i)}\le y\},
\qquad y<0.
\end{equation}

于是，对第 $j$ 条策略得分轨迹，其对应的 CPT 权重定义为
\begin{equation}
\psi_t^{(j)}
=
\begin{cases}
v_g\!\left(Y_t^{(j)}\right)\,
w_g'\!\left(\widehat p_t^g(Y_t^{(j)})\right),
& Y_t^{(j)}\ge 0,\\[1.2ex]
-\,v_l\!\left(-Y_t^{(j)}\right)\,
w_l'\!\left(\widehat p_t^l(Y_t^{(j)})\right),
& Y_t^{(j)}<0.
\end{cases}
\end{equation}

从而，第 $t$ 期的 CPT 策略梯度估计定义为
\begin{equation}
\label{eq:CPT-gradient-estimator}
\widehat{\nabla J_t}(\theta_{t-1},r_t)
=
\frac{1}{m_t}
\sum_{j=1}^{m_t}
\psi_t^{(j)}
\left(
\sum_{k=1}^{h}
\nabla_\theta
\log \pi_{\theta_{t-1}}
\bigl(
a_{t-k}^{(j)}\mid s_{t-k}^{(j)},c_t
\bigr)
\right).
\end{equation}

注意，式 \eqref{eq:CPT-gradient-estimator} 直接给出了第 $t$ 期的 CPT 策略梯度估计，本文不再单独引入经验 CPT 值 $\widehat J_t(\theta_{t-1},r_t)$ 的公式。

定义窗口条件下的目标梯度均值为
\begin{equation}
\nabla \bar J_t(\theta_{t-1},r_t)
=
\mathbb E
\!\left[
\widehat{\nabla J_t}(\theta_{t-1},r_t)
\mid H_t,r_t
\right].
\end{equation}

相应地，第 $t$ 期的梯度估计误差记为
\begin{equation}
\varepsilon_t
=
\widehat{\nabla J_t}(\theta_{t-1},r_t)
-
\nabla J_t(\theta_{t-1},r_t).
\end{equation}